% !PRINTED 57
\weberpage{57}

Es ist also $e$ von der Form $q^{\lambda-1}(q-1)$ mit einem positiven $\lambda$. Dass aber $\lambda$ nicht kleiner als $\varkappa$ sein kann, folgt aus (6), woraus zu sehen ist, dass $q^\lambda$ die höchste Potenz von $q$ ist, die in der Differenz

\[
g^{q^{\lambda-1}(q-1)} - 1
\]

aufgeht. Es ist also $c$ die kleinste positive Zahl, für die die Congruenz (8) erfüllt ist, und damit ist gleichbedeutend, dass die Zahlen

\weq{10}{1, g, g^2 \ldots g^{c-1}}

alle incongruent sind nach dem Modul $q^\varkappa$; $g$ kann also auch als \gesperrt{primitive Wurzel} von $q^\varkappa$ bezeichnet werden.

Die Anzahl der Glieder dieser Reihe (10) ist gleich $c$. Ebenso gross ist aber auch die Anzahl der nach dem Modul $q^\varkappa$ incongruenten, durch $q$ nicht theilbaren Zahlen $n$, und damit ist bewiesen:

\noindent\textbf{2.} \gesperrt{Wenn} $q$ \gesperrt{eine ungerade Primzahl,} $n$ \gesperrt{eine beliebige durch} $q$ \gesperrt{nicht theilbare Zahl,} $g$ \gesperrt{eine primitive Wurzel von} $q^2$ \gesperrt{und} $\varkappa$ \gesperrt{ein positiver Exponent ist, so lässt sich eine und nur eine Zahl} $\gamma$ \gesperrt{nach dem Modul} $c$ \gesperrt{bestimmen, die der Congruenz}

\[
n \equiv g^\gamma \pmod{q^\varkappa}
\]

\gesperrt{genügt.}

Die Zahl $c$ ist immer durch 2 theilbar, und wir heben also noch den besonderen Satz hervor, dass

\weq{11}{g^{\frac12 c} \equiv -1 \pmod{q^\varkappa}}

ist. Denn in dem durch $q^\varkappa$ theilbaren Producte

\[
g^c - 1 = (g^{\frac12 c} - 1)(g^{\frac12 c} + 1)
\]

ist der erste Factor nicht durch $q^\varkappa$ theilbar, und folglich muss der zweite Factor durch $q$ theilbar sein. Dann folgt aber, dass der erste Factor auch nicht durch $q$ theilbar sein kann, weil sonst auch die Summe der beiden Factoren, also auch $g$ selbst, durch $q$ theilbar sein müsste. Folglich muss der zweite Factor durch $q^\varkappa$ theilbar sein.

Hier besteht also vollständige Analogie mit den Sätzen über primitive Wurzeln und Indexsysteme, die wir im \S.\,136 des ersten Bandes kennen gelernt haben, und man kann also auch hier $\gamma$ als den Index von $n$ für den Modul $q^\varkappa$ bezeichnen.
